% Generated by scripts/assemble.py; corrections belong in source/corrections.json.
\examhead{82}{哈尔滨工业大学 2025 年数学分析试卷}{80}
% p081-b18; source PDF page 81
\blockanchor{p081-b18}
\question{1}判断题. 正确的给出证明, 错误的给出反例.\par
% p081-b19; source PDF page 81
\blockanchor{p081-b19}
(1) 设  $b_{n}=\frac{a_{1}+a_{2}+\cdots+a_{n}}{n}$ ，若数列  $\{b_{n}\}$  收敛，则  $\{a_{n}\}$  收敛.\par
% p081-b20; source PDF page 81
\blockanchor{p081-b20}
(2) 两个在 $x_0$ 附近无界的函数之积仍为无界函数.\par
% p081-b21; source PDF page 81
\blockanchor{p081-b21}
(3) 若 $y = f(x)$ 在点 $x_0$ 不连续, $u = g(y)$ 在点 $y_0 = f(x_0)$ 不连续, 则复合函数 $u = g(f(x))$ 在点 $x_0$ 不连续.\par
% p081-b22; source PDF page 81
\blockanchor{p081-b22}
(4) 一元函数的定积分, 若 $|f(x)|$ 可积, 则 $f(x)$ 可积.\par
% p081-b23; source PDF page 81
\blockanchor{p081-b23}
(5) 若二元函数 $f(x, y)$ 在点 $(x_0, y_0)$ 的两个二次极限都存在, 则 $f(x, y)$ 在 $(x_0, y_0)$ 的二重极限也存在.\par
% p081-b24; source PDF page 81
\blockanchor{p081-b24}
\question{2}解答如下问题:\par
% p081-b25; source PDF page 81
\blockanchor{p081-b25}
(1) 若函数 $f(x)$ 在 $(a, b)$ 上一致连续, 证明: $f(x)$ 在 $(a, b)$ 上有界.\par
% p081-b26; source PDF page 81
\blockanchor{p081-b26}
(2) 用两种方法证明函数 $y = \ln x$ 在 (0,1) 上不一致连续.\par
% p081-b27; source PDF page 81
\blockanchor{p081-b27}
\question{3}设函数 $f(x, y) = \begin{cases} (x^2 + y^2)^{\alpha} \sin \frac{1}{x^2 + y^2}, & x^2 + y^2 \neq 0; \\ 0, & x^2 + y^2 = 0. \end{cases}$\par
% p082-b00; source PDF page 82
\blockanchor{p082-b00}
(1) 求 $\alpha$ 的范围, 使函数 $f(x, y)$ 在 $(0, 0)$ 连续.\par
% p082-b01; source PDF page 82
\blockanchor{p082-b01}
(2) 求 $\alpha$ 的范围, 使函数 $f(x, y)$ 在 $(0, 0)$ 可微.\par
% p082-b02; source PDF page 82
\blockanchor{p082-b02}
(3) 求 $\alpha$ 的范围, 使函数 $f(x, y)$ 在 $(0,0)$ 的偏导数连续.\par
% p082-b03; source PDF page 82
\blockanchor{p082-b03}
\question{4}分别从微积分的角度, 用五种方法证明: $e^x \geq 1 + x$.\par
% p082-b04; source PDF page 82
\blockanchor{p082-b04}
\question{5}若对 $\forall x \in [a, b]$, 函数 $f(x)$ 满足 $f(x) \geq 0$, $f''(x) \leq 0$. 求证: $f(x) \leq \frac{2}{b - a} \int_{a}^{b} f(x) \mathrm{d}x$.\par
% p082-b05; source PDF page 82
\blockanchor{p082-b05}
\question{6}设函数 $f(x)$ 在 $(a, b)$ 上连续且无极大值点，$\{x_n\} \subset (a, b)$ 满足 $\lim_{n \to \infty} f(x_n) = \inf_{a < x < b} f(x)$.\par
% p082-b06; source PDF page 82
\blockanchor{p082-b06}
(1) 用肯定语言叙述 $\lim_{n\to \infty}x_n\neq \overline{x}$ 的 $\varepsilon -N$ 定义.\par
% p082-b07; source PDF page 82
\blockanchor{p082-b07}
(2) 若 $\lim_{n\to \infty}x_n = x_0$ ，是否有 $x_0\in (a,b)$ ？如果有,请加以证明，否则请举反例.\par
% p082-b08; source PDF page 82
\blockanchor{p082-b08}
(3) 证明 $\lim_{n\to \infty}x_n$ 存在.\par
% p082-b09; source PDF page 82
\par\Needspace{6\baselineskip}
\blockanchor{p082-b09}
\question{7}设正项级数 $\sum_{n=1}^{\infty} \frac{1}{p_n}$ 收敛. 记 $A_n = \sum_{k=1}^{n} p_k$, $B_n = \sum_{k=1}^{n} \frac{k^2}{A_k^2}$, $C_n = \sum_{k=1}^{n} \frac{1}{p_k}$, $S_n = \sum_{k=1}^{n} \frac{k}{A_k}$. 证明:\par
\nopagebreak[4]
% p082-b10; source PDF page 82
\blockanchor{p082-b10}
\displayformula{B _ {n} <   \frac {5}{p _ {1}} + 2 S _ {n} + C _ {n}. \tag {1}}
% p082-b11; source PDF page 82
\blockanchor{p082-b11}
\displayformula{S _ {n} <   \sqrt {B _ {n} C _ {n}}. \tag {2}}
% p082-b12; source PDF page 82
\blockanchor{p082-b12}
(3) $\sum_{n=1}^{\infty} \frac{n}{p_1 + p_2 + \cdots + p_n}$ 收敛.\par
% p082-b13; source PDF page 82
\blockanchor{p082-b13}
\question{8}证明:\par
% p082-b14; source PDF page 82
\blockanchor{p082-b14}
(1) $\sum_{n=1}^{\infty} x^{n} \ln x$ 在 $(0,1]$ 上不一致收敛.\par
% p082-b15; source PDF page 82
\blockanchor{p082-b15}
(2) $\sum_{n=1}^{\infty} x^n \frac{\ln x}{1 + |\ln(-\ln x)|}$ 在 (0,1) 上一致收敛.\par
% p082-b16; source PDF page 82
\blockanchor{p082-b16}
\question{9}设 $I(k, a) = \int_{0}^{+\infty} e^{-kx} \frac{\sin ax}{x} \, \mathrm{d}x$.\par
% p082-b17; source PDF page 82
\blockanchor{p082-b17}
(1) 证明: 对固定的 $k \in [0, +\infty)$, 积分 $I(k, a)$ 关于 $a$ 在 $|a| \geq \delta (\delta > 0)$ 上一致收敛. 对固定的 $a \neq 0$, 积分 $I(k, a)$ 关于 $k$ 在 $[0, +\infty)$ 上一致收敛.\par
% p082-b18; source PDF page 82
\blockanchor{p082-b18}
(2) 计算积分 $I = \int_{0}^{+\infty} \frac{\sin x}{x} \, \mathrm{d}x$ 的值.\par
% p082-b19; source PDF page 82
\blockanchor{p082-b19}
\question{10}计算:\par
% p082-b20; source PDF page 82
\blockanchor{p082-b20}
(1) $\oint_{C} \frac{y}{x^{2} + 4y^{2}} \mathrm{d}x - \frac{x}{x^{2} + 4y^{2}} \mathrm{d}y,$ 其中 $C$ 是圆周 $x^{2} + y^{2} = 1$, 取逆时针方向.\par
% p082-b21; source PDF page 82
\blockanchor{p082-b21}
(2) $\iint_{S} (y - z) \, \mathrm{d}y \, \mathrm{d}z + (z - x) \, \mathrm{d}z \, \mathrm{d}x + (x - y) \, \mathrm{d}x \, \mathrm{d}y$, 其中 $S$ 是圆锥 $x^2 + y^2 = z^2 (0 \leq z \leq 2)$ 的外侧.\par
% p082-b22; source PDF page 82
\blockanchor{p082-b22}
(3) $\oint_{C} y^{3} \mathrm{d}x + z^{3} \mathrm{d}y + x^{3} \mathrm{d}z$, 其中 $C$ 为球面 $x^{2} + y^{2} + z^{2} = 4$ 与平面 $x + y + z = 0$ 的交线, 从 $z$ 轴正向看为逆时针方向.\par
